\documentclass[12pt]{article}
\usepackage[utf8]{inputenc}
\usepackage[T1]{fontenc}
\usepackage{hyperref}
\title{Megleno-Romanian Language Resources for K-12 Curriculum Development}
\author{Compiled by Claude Sonnet 5.5 (claude-sonnet-5-5) for Andr\'as Kornai}
\date{}
\begin{document}
\maketitle
\section{Introduction}
Megleno-Romanian (ISO 639-3 ruq) is an endangered Romance language of Greece and North Macedonia. This dossier is based on a single search and is short. Pages were not opened.

\section{Primary Educational Resources}
No school textbooks, primers or programmes were found.

\section{Linguistic and Reference Materials}
\begin{itemize}
\item Bibliography entry in Diacronia.ro that cites Capidan's \emph{Meglenorom\^anii}, vol.\ III (dictionary, 1935) and Candrea's glossary in \emph{Grai \c{s}i suflet}. \url{https://www.diacronia.ro/indexing/details/A29397/pdf}
\item Another Diacronia.ro record from the same results. \url{https://www.diacronia.ro/indexing/details/A15309/pdf}
\item Glottolog entry for Megleno-Romanian. \url{https://glottolog.org/resource/languoid/id/megl1237}
\item Omniglot page on Megleno-Romanian (vl\u{a}\c{s}\'e\c{s}te). \url{https://www.omniglot.com/writing/meglenoromanian.htm}
\end{itemize}

\section{Religious and Community Resources}
None found.

\section{Digital and Audio Resources}
\begin{itemize}
\item Glosbe section for ruq. Coverage is uneven and it is crowd-sourced. \url{https://glosbe.com/lib/ruq}
\item Verbix verb conjugation page. \url{https://docs.verbix.com/Languages/Meglenoromanian}
\end{itemize}

\section{Curriculum Development Context}
Little usable material was found. A curriculum would need to start from the early 20th-century scholarly dictionaries and texts, plus fieldwork with speakers.
\end{document}